% Source-aligned English for printed p.132.
\SourcePage{132}

\LayerTarget{tohoku:example:1.7:b:p2}
\noindent commutation relations, one obtains a product of categories of the
kind considered in No.~1.

\LayerTarget{tohoku:example:1.7:c}
\noindent c) Take $I$ to consist of two elements, $a$ and $b$, with a single
arrow having origin $a$ and endpoint $b$. This gives the category of
morphisms $u:A\longrightarrow B$ between objects of $C$. Introducing
commutation relations would amount to restricting to $A,B,u$ annihilated by
specified integers.

\LayerTarget{tohoku:example:1.7:d}
\noindent d) \textit{Categories of functors.} Let $C'$ be another category,
assumed to be a set. The covariant functors from $C'$ to $C$ then form a
category, with natural transformations as morphisms (cf. 1, 1). This category
may be interpreted as a category $C^\Sigma$: take $I=C'$, let the arrows with
origin $A'$ and endpoint $B'$ be, by definition, the elements of
$\mathrm{Hom}(A',B')$, and take as commutation relations those expressing the
two functor axioms. If $C'$ is moreover an additive category, the
\textit{additive} functors from $C'$ to $C$ may likewise be interpreted as a
category $C^\Sigma$ (by adjoining the required commutation relations).

\LayerTarget{tohoku:example:1.7:e}
\noindent e) \textit{Complexes with values in $C$.} Let $I=\mathbf{Z}$, the
set of integers, and let the arrows be $(d_n)_{n\in\mathbf{Z}}$, where $d_n$
has origin $n$ and endpoint $n+1$; the commutation relations are
$d_{n+1}d_n=0$. Relations of the form $e_n=0$ may also be adjoined if one
wishes to restrict to complexes in positive degrees, or in negative degrees.
Bicomplexes and so forth are obtained analogously.

\LayerTarget{tohoku:example:1.7:f}
\noindent f) \textit{The category $C^G$ ($G$ a group).} Let $G$ be a group and
$C$ a category, not necessarily additive. An \textit{object in $C$ with
operator group $G$} is a pair $(A,r)$ consisting of an object $A\in C$ and a
representation $r$ of $G$ in the automorphism group of $A$. If $(A',r')$ is a
second such pair, a morphism from the first to the second is a morphism from
$A$ to $A'$ commuting with the operations of $G$. The class $C^G$ of objects
of $C$ with operator group $G$ thereby becomes a category. It may be
interpreted as a class $C^\Sigma$, where $\Sigma=\Sigma(G)$ is the following
scheme with relations: its vertex set consists of a single element $i_0$, its
set of arrows is $G$, and the commutation relations are $(s)(t)=(st)$, where
the left-hand side denotes a composite arrow, and $(e)=e_{i_0}$, where $e$
denotes the identity element of $G$. In particular, if $C$ is additive, then
so is $C^G$; in that case this construction is contained in the next one, by
considering the group algebra of $G$.

\LayerTarget{tohoku:example:1.7:g}
\noindent g) \textit{The category $C^U$ ($U$ a ring with identity).} Consider
the additive category formed by pairs $(A,r)$ consisting of an object $A$ of
$C$ and a unital representation of $U$ in the ring $\mathrm{Hom}(A,A)$, with
morphisms defined in the evident way. As above, it is interpreted as a
category $C^{\Sigma(U)}$, where $\Sigma(U)$ is the scheme with relations
having one vertex, $U$ as its set of arrows, and commutation relations that we
shall not write down.

\LayerTarget{tohoku:example:1.7:h:p1}
\noindent h) \textit{Inductive systems and projective systems.} Take as the
set of vertices a \textit{preordered} set $I$, and as arrows the pairs $(i,j)$
of vertices with $i\leq j$, the origin and endpoint of $(i,j)$ being $i$ and
$j$, respectively. The commutation relations are
$(i,j)(j,k)=(i,k)$ and $(i,i)=e_i$. The

\TohokuContinuation{133}
